\documentclass[11pt,a4paper]{article}
% Shared preamble for the English probability and statistics documents.
% Based on the original course preamble. Compile with pdfLaTeX.
\usepackage[top=1.5cm,bottom=2cm,left=2cm,right=2cm]{geometry}
\usepackage[T1]{fontenc}
\usepackage{lmodern}
\usepackage[utf8]{inputenc}
\usepackage[english]{babel}

\usepackage{amsmath}
\usepackage{amssymb}
\usepackage{amsthm}
\usepackage{mathrsfs}
\usepackage{mathtools}
\usepackage{etoolbox}
\usepackage{tikz}
\usetikzlibrary{shapes.multipart,arrows.meta}
\usepackage{xcolor}
\usepackage{enumerate}
\usepackage[shortlabels]{enumitem}
\usepackage{hyperref}
\usepackage{version}
\usepackage{bbold}
\usepackage{fontawesome5}
\usepackage{fancyhdr}
\lfoot{}
\usepackage{natbib}
\usepackage{graphicx}
\usepackage{pdfpages}
\usepackage{authblk}
    
\usepackage{blindtext}
\usepackage{multicol}
\usepackage{collcell}
\usepackage{wrapfig}
\usepackage{blkarray}
\usepackage{cancel}
\usepackage{appendix}
\usepackage{float}
\makeatletter
\input{tex/assets/slashbox.sty}
\makeatother
\hypersetup{unicode,colorlinks,citecolor=red,filecolor=black,
  linkcolor=red,urlcolor=cyan}
\setlength{\columnsep}{1cm}
\pagenumbering{gobble}

\def\numtext#1{%
  \ifcase#1\or one\or two\or three\or four\or five\or six\or seven\or eight\or nine\or ten\or
  eleven\or twelve\or thirteen\or fourteen\or fifteen\or sixteen\or seventeen\or eighteen\or nineteen\or twenty\or Lots\fi}

\newcommand{\Byte}[1]{%
  \def\nodecontents{}%
  \foreach[count=\l] \k in {#1} {\xappto{\nodecontents}{\noexpand\nodepart{\numtext{\l}}\k}}%
  \node[rectangle split,rectangle split parts=8,rectangle split horizontal,draw] at (2,2){\nodecontents};}

\newcommand{\D}{{\mathrm{d}}}
\newcommand{\cE}{{\mathcal{E}}}
\newcommand{\cG}{{\mathcal{G}}}
\newcommand{\cX}{{\mathcal{X}}}
\newcommand{\cY}{{\mathcal{Y}}}
\newcommand{\cB}{{\mathcal{B}}}
\newcommand{\cF}{{\mathcal{F}}}
\newcommand{\cO}{{\mathcal{O}}}
\newcommand{\cS}{{\mathcal{S}}}
\newcommand{\cT}{{\mathcal{T}}}
\newcommand{\given}{\, \big| \,}
\newcommand{\cL}{{\mathscr{L}}}
\newcommand{\cD}{{\mathcal{D}}}
\newcommand{\cA}{{\mathscr{A}}}
\newcommand{\cP}{{\mathscr{P}}}
\newcommand{\cR}{{\mathscr{R}}}
\newcommand{\A}{{\mathcal A}}
\newcommand{\one}{\mathbb{1}}
\newcommand{\mathbbm}[1]{\mathbb{#1}}

\providecommand{\C}{}
\renewcommand{\C}{\mathbb C}
\newcommand{\R}{\mathbb R}
\newcommand{\Q}{\mathbb Q}
\newcommand{\N}{\mathbb N}
\newcommand{\T}{\mathbb T}
\newcommand{\F}{\mathbb F}
\renewcommand{\P}{\mathbb P}
\newcommand{\Z}{\mathbb Z}
\newcommand{\E}{\mathbb {E}}
\newcommand{\Pbb}[1]{\P\left( #1\right)}
\newcommand{\Ebb}[1]{\E\left[ #1\right]}

\newcommand{\lb}{\left (}
\newcommand{\rb}{\right )}
\newcommand{\lbb}{\left [}
\newcommand{\rbb}{\right ]}
\newcommand{\labs}{\left |}
\newcommand{\rabs}{\right |}
\newcommand{\lbrb}[1]{\lb #1 \rb}
\newcommand{\lbbrbb}[1]{\lbb#1\rbb}
\newcommand{\labsrabs}[1]{\labs#1\rabs}
\newcommand{\lbbrb}[1]{\lbb#1\rb}
\newcommand{\lbrbb}[1]{\lb#1\rbb}
\newcommand{\langlerangle}[1]{\langle#1\rangle}
\newcommand{\lbcurly}{\{}
\newcommand{\rbcurly}{\}}
\newcommand{\curly}[1]{\lbcurly#1\rbcurly}
\newcommand{\abs}{\labsrabs}
\newcommand{\ind}[1]{1_{#1}}
\newcommand{\indep}{\perp \!\!\! \perp}
\renewcommand{\d}{\, \mathrm{d}}

\newtheorem{theorem}{Theorem}
\newtheorem{problemCnt}{Lemma}
\theoremstyle{definition}
\newtheorem{problem}[problemCnt]{Problem}
\newtheorem{problembook}[problemCnt]{\faBook \,  Problem}
\newtheorem{problempencil}[problemCnt]{\faPencil* \,  Problem}
\newtheorem{problemstar}[problemCnt]{\faStar\,  Problem}
\newtheorem{problemstarpencil}[problemCnt]{\faStar \, \faPencil* \,   Problem}
\newtheorem{probleminteresting}[problemCnt]{\faSmile[regular] \,   Problem}
\newtheorem{solution}{Solution}

\newcommand{\fwcell}[1]{\makebox[\arraystretch\normalbaselineskip]{$#1$}}
\newenvironment{enumeratei}{\begin{enumerate}[(i)]}{\end{enumerate}}
\newenvironment{enumeratea}{\begin{enumerate}[(a)]}{\end{enumerate}}
\newenvironment{enumeraten}{\begin{enumerate}[\upshape 1.]}{\end{enumerate}}

% Compatibility helpers needed by a few files in the combined archive.
\AtBeginDocument{\noindent{\footnotesize English translation prepared with ChatGPT (OpenAI, GPT-6), 26 October 2026, for translation of the Bulgarian text into English and LaTeX formatting.
The Bulgarian originals remain the reference; translation may contain errors.}\par\smallskip}
\newcommand{\MissingGraphic}[1]{%
  \begin{center}\fbox{\parbox{0.82\linewidth}{\centering Missing source image: \texttt{#1}}}\end{center}}



\pagestyle{fancy}\fancyhf{}\lhead{Applied Mathematics: Test 2 solutions}\rhead{8 June 2024}
\setlength{\headheight}{15pt}\pagenumbering{arabic}\cfoot{\thepage}
\begin{document}
\noindent\emph{English transcription of the original handwritten solution notes. Cross-references and the level of detail of the original are retained.}

\section*{Problems 1 and 2}
See Problems 1 and 2 of the Computer Science Test 2, 2024.

\section*{Problem 3}
See the exercise classes.

\section*{Problem 4}
Let $X,Y\sim\operatorname{Exp}(1)$ be independent and identically distributed.
Recall that $\P(\operatorname{Exp}(\lambda)>t)=e^{-\lambda t}$.
For $t>0$,
\begin{align*}
\P\left(\sqrt{X/Y}>t\right)
&=\P(X>Yt^2)\\
&=\int_0^\infty\P(X>yt^2)f_Y(y)\d y\\
&=\int_0^\infty e^{-yt^2}e^{-y}\d y\\
&=\frac1{1+t^2}\int_0^\infty e^{-y(1+t^2)}\d[y(1+t^2)]
=\frac1{1+t^2},
\end{align*}
using $\int_0^\infty e^{-x}\d x=1$. Therefore,
\[
\E\sqrt{X/Y}=\int_0^\infty\P\left(\sqrt{X/Y}>t\right)\d t
=\int_0^\infty\frac{\d t}{1+t^2}=\frac\pi2.
\]
Alternatively, use $F_{\sqrt{X/Y}}(t)=t^2/(1+t^2)$ and obtain the density by differentiating, $f=F'$, and so on.

\subsection*{Second method}
\begin{align*}
\E\sqrt{X/Y}
&=\int_0^\infty\int_0^\infty\sqrt{x/y}\,f_{X,Y}(x,y)\d x\d y\\
&=\int_0^\infty\int_0^\infty\sqrt{x/y}\,e^{-x}e^{-y}\d x\d y\\
&=\left(\int_0^\infty e^{-x}\sqrt x\d x\right)
\left(\int_0^\infty e^{-y}\frac1{\sqrt y}\d y\right)\\
&=\Gamma(3/2)\Gamma(1/2)=\frac12\Gamma(1/2)^2,
\end{align*}
where $\Gamma(z+1)=z\Gamma(z)$. By the first method,
\[
\frac12\Gamma(1/2)^2=\frac\pi2,
\qquad\text{hence}\qquad\Gamma(1/2)=\sqrt\pi.
\]

\section*{Problem 5}
See Problem 4 of the Software Engineering examination, February 2024.

\section*{Problem 6}
This is the proof of a special case of Cochran's theorem, divided into individual steps.

\section*{Problem 7}
\[
e^{-n}\left(\frac{n^0}{0!}+\frac{n^1}{1!}+\cdots+\frac{n^n}{n!}\right)
=\sum_{k=0}^n\P(\operatorname{Poi}(n)=k)
=\P(\operatorname{Poi}(n)\le n).
\]
As seen earlier in the course, if $X_n\sim\operatorname{Poi}(n)$, then
\[
X_n\overset d=Y_1+\cdots+Y_n,
\qquad Y_1,\ldots,Y_n\ \text{i.i.d.}\ \operatorname{Poi}(1),
\qquad\E Y_i=D(Y_i)=1.
\]
Thus, by the central limit theorem,
\begin{align*}
\P(\operatorname{Poi}(n)\le n)
&=\P(Y_1+\cdots+Y_n\le n)\\
&=\P\left(\frac{Y_1+\cdots+Y_n-n}{\sqrt n}
\le\frac{n-n}{\sqrt n}\right)\\
&\longrightarrow\P(N(0,1)\le0)=\frac12.
\end{align*}
\end{document}
